\section{Automations：自动化调度}
%% ============================================================

\subsection{从 Skill 到 Automation}

当一个工作流已经稳定可靠，就可以用 Automation 让 Codex 在后台按计划自动执行。

\begin{importantbox}{Skills vs Automations}
\textbf{Skills 定义方法，Automations 定义调度。}如果一个工作流仍然需要大量人工干预，应当先将其封装为 skill。只有当它足够可预测时，automation 才能真正发挥乘数效应。
\end{importantbox}

\subsection{适用场景}

好的 automation 候选包括：

\begin{itemize}
    \item 汇总近期 commits
    \item 扫描潜在 bug
    \item 起草 release notes
    \item 检查 CI 失败
    \item 生成 standup 摘要
    \item 按计划运行重复性分析
\end{itemize}

在 Codex App 的 Automations 标签页中，可以选择：
\begin{itemize}
    \item 运行的项目
    \item 执行的 prompt（可调用 skills）
    \item 运行的节奏/频率
    \item 执行环境（dedicated git worktree 或本地环境）
\end{itemize}

\begin{warningbox}{先手动验证，再自动化}
在工作流手动运行已经可靠之前，不要将其转为 automation。过早自动化不可靠的任务只会放大错误。
\end{warningbox}

\begin{knowledgebox}{用 Automation 做反思和维护}
Automation 不仅仅用于执行任务。还可以用来审查近期 session、汇总重复性摩擦、改进 prompt 和工作流设置。
\end{knowledgebox}

\subsection{本章小结}

Automation 是将稳定 skill 按计划自动执行的机制。先确保手动运行可靠，再自动化。可用于执行，也可用于反思和维护。

%% ============================================================
